\documentclass[12pt, a4paper,oneside, UTF8]{ctexbook}
\usepackage[dvipsnames]{xcolor}
\usepackage{amsmath}   % 数学公式
\usepackage{booktabs} % 制作三线表
\usepackage{ntheorem}   % 定理环境
%\usepackage{amsthm}    % 定理环境
\usepackage{amssymb}   % 更多公式符号
\usepackage{graphicx}  % 插图
\usepackage{mathrsfs}  % 数学字体
\usepackage{enumitem}  % 列表
\usepackage{geometry}  % 页面调整
\usepackage{lipsum}    % 自动生成一段文本
\usepackage[strict]{changepage} % 提供一个 adjustwidth 环境
\usepackage{framed} % 实现方框效果
\usepackage{listings}% 代码环境
\usepackage{tikz}
\usepackage{unicode-math}
\usepackage[colorlinks,linkcolor=black]{hyperref}
\usepackage{minted}% 代码环境
\usemintedstyle{xcode}
\graphicspath{ {flg/},{../flg/}, {config/}, {../config/} }  % 配置图形文件检索目录
\linespread{1.5} % 行高

% 页码设置
\geometry{top=25.4mm,bottom=25.4mm,left=20mm,right=20mm,headheight=2.17cm,headsep=4mm,footskip=12mm}

% 设置列表环境的上下间距
%\setenumerate[1]{itemsep=5pt,partopsep=0pt,parsep=\parskip,topsep=5pt}
%\setitemize[1]{itemsep=5pt,partopsep=0pt,parsep=\parskip,topsep=5pt}
%\setdescription{itemsep=5pt,partopsep=0pt,parsep=\parskip,topsep=5pt}

% 自定义颜色
\definecolor{formalshade}{rgb}{0.95,0.95,1} % 文本框颜色
\definecolor{codeframeshade}{RGB}{231,239,250}
% 设置章形式
\usepackage{titlesec, titletoc}
\linespread{1.2}                
\usepackage{fancyhdr}
\fancyhf{}
\renewcommand{\headrule}{\color{black}\hrule width\textwidth}
\pagestyle{fancy}
\renewcommand{\headrulewidth}{1pt}
\fancypagestyle{plain}{\renewcommand{\headrulewidth}{0pt}\fancyhf{}\renewcommand{\headrule}{}}

\fancyhead[c]{\color{red}\kaishu\rightmark}
\fancyfoot[c]{\color{red}\small\thepage}

\titleformat{\chapter}[display]{\Large}
{\color{black}\filleft
	\parbox{1cm}{\vbox to 1.5cm{\vfill\hbox to 4cm{\hfill\Huge \bfseries \color{black}{Chapter} \thechapter \hfill}}}}
{1ex}
{\color{black} \titlerule[2pt]\large\bfseries \filright \vspace*{1em}}
[\vspace*{1em} {\titlerule[2pt]}]

% 在导言区进行代码样式设置
\lstset{
	language=Java, % 设置语言
	basicstyle=\ttfamily, % 设置字体族
	breaklines=true, % 自动换行
	keywordstyle=\bfseries\color{NavyBlue}, % 设置关键字为粗体，颜色为 NavyBlue
	morekeywords={}, % 设置更多的关键字，用逗号分隔
	emph={self}, % 指定强调词，如果有多个，用逗号隔开
	emphstyle=\bfseries\color{Rhodamine}, % 强调词样式设置
	commentstyle=\itshape\color{black!50!white}, % 设置注释样式，斜体，浅灰色
	stringstyle=\bfseries\color{PineGreen!90!black}, % 设置字符串样式
	columns=flexible,
	numbers=left, % 显示行号在左边
	numbersep=1em, % 设置行号的具体位置
	numberstyle=\footnotesize, % 缩小行号
	stepnumber=5,% 行号设置
	% frame=single, % 边框
	rulecolor=\color{gray!90},      % 代码框颜色
	aboveskip=.25em, 			% 代码块边框
	framextopmargin=0pt, framexbottommargin=0pt,
	framesep=2em, % 设置代码与边框的距离
	xleftmargin=1cm,
	xrightmargin=1cm,
	tabsize=3
}
% 定理环境
%设置定理环境的样式：plain(默认样式)、difinition、remark
%\theoremstyle{definition}
\theorembodyfont{\kaishu}
%\newtheorem{环境名}{标题}[排序单位]
\newtheorem{definition}{\indent 定义}[section]
\newtheorem{lemma}{\indent 引理}[section]    
\newtheorem{theorem}{\indent 定理}[section]
\newtheorem{corollary}{\indent 推论}[section]
% 引理 定理 推论 准则 共用一个编号计数
%\newtheorem{thm}[lemma]{\indent 定理}
%\newtheorem{corollary}[lemma]{\indent 推论}
%\newtheorem{criterion}[lemma]{\indent 准则}
%\newtheorem{proposition}{\indent 命题}[section]

\newtheorem{example}{\indent \color{SeaGreen}{例}} % 绿色文字的例 ， 不需要可更改 \color{SeaGreen}{} 
\newtheorem*{remark}{\indent 注}

% 定义中文的 证明 和 解 的环境：
\newenvironment{proof}{\par\textbf{证明}\;}{
	\vspace{-2.5em}
	\begin{flushright}
		$\square$
	\end{flushright}\par}
\newenvironment{solution}{\par{\textbf{解}}\;}{\par}


% ↓↓↓↓↓↓↓↓↓↓↓↓↓↓↓↓↓ 以下是自定义的命令  ↓↓↓↓↓↓↓↓↓↓↓↓↓↓↓↓

% 用于调整表格的高度  使用 \hline\xrowht{25pt}
\newcommand{\xrowht}[2][0]{\addstackgap[.5\dimexpr#2\relax]{\vphantom{#1}}}

% 表格环境内长内容换行
\newcommand{\tabincell}[2]{\begin{tabular}{@{}#1@{}}#2\end{tabular}}

% 使用\linespread{1.5} 之后 cases 环境的行高也会改变，重新定义一个 ca 环境可以自动控制 cases 环境行高
\newenvironment{ca}[1][1]{\linespread{#1} \selectfont \begin{cases}}{\end{cases}}
% 和上面一样
\newenvironment{vx}[1][1]{\linespread{#1} \selectfont \begin{vmatrix}}{\end{vmatrix}}

\def\d{\textup{d}} % 直立体 d 用于微分符号 dx
\def\R{\mathbb{R}} % 实数域
\newcommand{\bs}[1]{\boldsymbol{#1}}    % 加粗，常用于向量
\newcommand{\ora}[1]{\overrightarrow{#1}} % 向量

% 数学 平行 符号
\newcommand{\pll}{\kern 0.56em/\kern -0.8em /\kern 0.56em}

% 用于空行\myspace{1} 表示空一行 填 2 表示空两行  
\newcommand{\myspace}[1]{\par\vspace{#1\baselineskip}}

\newcommand{\HRule}{\rule{\linewidth}{0.5mm}}%\rule[水平高度]{长度}{粗细}


% ------------------******-------------------
% 注意行末需要把空格注释掉，不然画出来的方框会有空白竖线
\newenvironment{formal}
{%
	\def\FrameCommand
	{%
		\hspace{1pt}%
		{\color{blue}\vrule width 2pt}%
		{\color{formalshade}\vrule width 4pt}%
		\colorbox{formalshade}%
	}%
	\MakeFramed{\advance\hsize-\width\FrameRestore}%
	\noindent\hspace{-4.55pt}% disable indenting first paragraph
	\begin{adjustwidth}{}{7pt}%
		\vspace{2pt}\vspace{2pt}%
	}
	{%
		\vspace{2pt}\end{adjustwidth}\endMakeFramed%
}

